\subsection{Ones-Counter System}
\label{sec:case-ones}

For readers less familiar with T3SD: a \emph{reference system} $Z_{\text{spec}}$ is a complete discrete-system table (states, inputs, outputs, next-state, readout) that expresses required behavior. A \emph{buildable system} $Z_{\text{impl}}$ is another such table---often with richer internal state---that is intended to realize the same behavior. Classical checking asks for a homomorphism (projection maps) from $Z_{\text{impl}}$ onto $Z_{\text{spec}}$. Here we instead compile $\Phi_{\mathrm{dyn}}$ from $Z_{\text{spec}}$ and check that property set on each lifecycle iteration of $Z_{\text{impl}}$.

\paragraph{Reference system $Z_{\mathrm{ones}}$.}
Behavior in plain language: start at count $0$; on input $1$ increment the count; on input $0$ leave it unchanged; the output at every tick is the current count.

\begingroup
\renewcommand{\arraystretch}{1.25}
\footnotesize
\begin{center}
\begin{tabularx}{0.78\linewidth}{@{}>{\raggedright\arraybackslash}p{2.8cm}Y@{}}
    \rowcolor{tableheader}
    \headingfont\bfseries Component & \headingfont\bfseries $Z_{\mathrm{ones}}$ \\
    \addlinespace[2pt]
    $S,I,O$ & $\mathbb{N}$, $\{0,1\}$, $\mathbb{N}$ \\
    \addlinespace[2pt]
    $s_0$ & $0$ \\
    \addlinespace[2pt]
    $\delta(n,1)$ & $n+1$ \\
    \addlinespace[2pt]
    $\delta(n,0)$ & $n$ \\
    \addlinespace[2pt]
    $\delta(n,\mathrm{none})$ & $n$ \\
    \addlinespace[2pt]
    $\rho(n)$ & $n$ \\
\end{tabularx}
\end{center}
\captionof{table}{Ones-counter reference system $Z_{\mathrm{ones}}$.}
\label{tab:z-ones}
\endgroup
\vspace{4pt}

\paragraph{Compiled $\Phi_{\mathrm{dyn}}$.}
For every $n\in\mathbb{N}$:
\begin{align*}
  &\text{(readout)} && G\bigl(\mathit{state}(n) \rightarrow \mathit{output}(n)\bigr), \\
  &\text{(guided $1$)} && G\bigl((\mathit{state}(n)\wedge\mathit{input}(1)) \rightarrow X(\mathit{state}(n{+}1))\bigr), \\
  &\text{(guided $0$)} && G\bigl((\mathit{state}(n)\wedge\mathit{input}(0)) \rightarrow X(\mathit{state}(n))\bigr), \\
  &\text{(autonomous)} && G\bigl((\mathit{state}(n)\wedge\mathit{input}(\mathrm{none})) \rightarrow X(\mathit{state}(n))\bigr).
\end{align*}
Denote this family by $\Phi_{Z_{\mathrm{ones}}}$. Hom-relative checking uses $\mathit{specState}(n)$ meaning ``the buildable's projected count equals $n$.''

\paragraph{Lifecycle iterations.}
All three share inputs $\{0,1\}$ and output the ones-count. They differ in internal bookkeeping. Each is checked against the \emph{same} $\Phi_{Z_{\mathrm{ones}}}$. This ladder is pedagogical: elaborations remain mild, and $h_S$ stays a projection onto the ones component.

\textit{v1 --- prototype $Z_{\mathrm{ones}}^{\mathrm{v1}}$.}
Identical to Table~\ref{tab:z-ones}. The homomorphism is the identity: $h_S=h_I=h_O=\mathrm{id}$.

\textit{v2 --- instrumented $Z_{\mathrm{ones}}^{\mathrm{v2}}$.}
State is a pair $(\mathrm{ones},\mathrm{zeros})$. On input $1$ increment ones; on input $0$ increment zeros; output is ones only.

\begingroup
\renewcommand{\arraystretch}{1.2}
\footnotesize
\begin{center}
\begin{tabularx}{0.78\linewidth}{@{}>{\raggedright\arraybackslash}p{2.8cm}Y@{}}
    \rowcolor{tableheader}
    \headingfont\bfseries Component & \headingfont\bfseries $Z_{\mathrm{ones}}^{\mathrm{v2}}$ \\
    \addlinespace[2pt]
    $S$ & $\mathbb{N}\times\mathbb{N}$ (ones, zeros) \\
    \addlinespace[2pt]
    $s_0$ & $(0,0)$ \\
    \addlinespace[2pt]
    $\delta((o,z),1)$ & $(o{+}1,\,z)$ \\
    \addlinespace[2pt]
    $\delta((o,z),0)$ & $(o,\,z{+}1)$ \\
    \addlinespace[2pt]
    $\rho(o,z)$ & $o$ \\
    \addlinespace[2pt]
    Maps to $Z_{\mathrm{ones}}$ & $h_S(o,z)=o$,\ $h_I=\mathrm{id}$,\ $h_O=\mathrm{id}$ \\
\end{tabularx}
\end{center}
\captionof{table}{Instrumented ones-counter (v2): tracks zeros as well, but only ones are required by $\Phi_{Z_{\mathrm{ones}}}$.}
\label{tab:z-ones-med}
\endgroup
\vspace{4pt}

\textit{v3 --- richer logging $Z_{\mathrm{ones}}^{\mathrm{v3}}$.}
State is a 4-tuple $(\mathrm{ones},\mathrm{ticks},\mathrm{lastBit},\mathrm{shadow})$. Ticks always advance; lastBit stores the previous input; shadow is free instrumentation (here updated by adding the current ones-count). Output remains ones.

\begingroup
\renewcommand{\arraystretch}{1.2}
\footnotesize
\begin{center}
\begin{tabularx}{0.78\linewidth}{@{}>{\raggedright\arraybackslash}p{2.8cm}Y@{}}
    \rowcolor{tableheader}
    \headingfont\bfseries Component & \headingfont\bfseries $Z_{\mathrm{ones}}^{\mathrm{v3}}$ \\
    \addlinespace[2pt]
    $S$ & $\mathbb{N}\times\mathbb{N}\times\{0,1\}\times\mathbb{N}$ \\
    \addlinespace[2pt]
    $s_0$ & $(0,0,0,0)$ \\
    \addlinespace[2pt]
    $\delta((o,t,\ell,s),1)$ & $(o{+}1,\,t{+}1,\,1,\,s{+}o)$ \\
    \addlinespace[2pt]
    $\delta((o,t,\ell,s),0)$ & $(o,\,t{+}1,\,0,\,s{+}o)$ \\
    \addlinespace[2pt]
    $\rho(o,t,\ell,s)$ & $o$ \\
    \addlinespace[2pt]
    Maps to $Z_{\mathrm{ones}}$ & $h_S(o,t,\ell,s)=o$,\ $h_I=\mathrm{id}$,\ $h_O=\mathrm{id}$ \\
\end{tabularx}
\end{center}
\captionof{table}{Richer ones-counter (v3): instrumentation; only the ones component is constrained by $\Phi_{Z_{\mathrm{ones}}}$.}
\label{tab:z-ones-cplx}
\endgroup
\vspace{4pt}

\paragraph{Workflow contrast (after v3).}
\textit{Assertional:} $\Phi_{Z_{\mathrm{ones}}}$ is unchanged from v1; verifying v2 or v3 means re-checking that same fragment on the new tables (no new derivation of maps to answer the implementability question).
\textit{Classical:} for $Z_{\mathrm{ones}}^{\mathrm{v3}}$, take $h_S(o,t,\ell,s)=o$, $h_I=\mathrm{id}$, $h_O=\mathrm{id}$. Surjectivity is immediate. Transition preservation for input $1$:
\[
  h_S\bigl(\delta_{\mathrm{v3}}((o,t,\ell,s),1)\bigr)
  = o{+}1
  = \delta_{\mathrm{ones}}\bigl(h_S(o,t,\ell,s),\,h_I(1)\bigr),
\]
and likewise for $0$ and $\mathrm{none}$; readout $\rho_{\mathrm{v3}}=o=\rho_{\mathrm{ones}}\circ h_S$. The witness remains a projection---honest that this pedagogical ladder does not yet force a hard $h_S$---but each elaboration still requires a human to confirm that the chosen projection still preserves $\delta$ and $\rho$. By the bi-implication, $Z_{\mathrm{ones}}^{\mathrm{v3}}\models_{\mathrm{dyn}}\Phi_{Z_{\mathrm{ones}}}$.

\paragraph{What checking means.}
For $Z_{\mathrm{ones}}^{\mathrm{v2}}$ and $Z_{\mathrm{ones}}^{\mathrm{v3}}$, $\mathit{specState}(n)$ holds when $h_S$ of the current buildable state equals $n$. Extra components may evolve freely so long as they do not disturb that projection. Revising the shadow update does not change $\Phi_{Z_{\mathrm{ones}}}$; one re-checks the same property set.

\subsection{String Recognizer for \texorpdfstring{$01110$}{01110}}
\label{sec:case-pattern}

\paragraph{Reference system $Z_{01110}$.}
Behavior: walk a small overlap-aware automaton that tracks how much of the bit-string $01110$ has been matched so far. Output $1$ only in the ``just matched'' state; otherwise output $0$. After a match, input $1$ leaves progress $2$ (prefix $01$), not a full reset---so the automaton agrees with longest-prefix-suffix progress on the recent window.

\begin{figure}[t]
\centering
\begin{tikzpicture}[
  >=Stealth, node distance=1.35cm and 1.1cm,
  every node/.style={font=\scriptsize},
  state/.style={circle, draw, minimum size=0.72cm, inner sep=0pt},
  acc/.style={state, double}
]
  \node[state, initial, initial where=left, initial text=] (q0) {$0$};
  \node[state, right=of q0] (q1) {$1$};
  \node[state, right=of q1] (q2) {$2$};
  \node[state, right=of q2] (q3) {$3$};
  \node[state, right=of q3] (q4) {$4$};
  \node[acc, right=of q4] (q5) {$5$};
  \path[->]
    (q0) edge node[above] {$0$} (q1)
    (q1) edge node[above] {$1$} (q2)
    (q2) edge node[above] {$1$} (q3)
    (q3) edge node[above] {$1$} (q4)
    (q4) edge node[above] {$0$} (q5);
  \path[->]
    (q0) edge[loop above] node {$1$} (q0)
    (q1) edge[loop below] node {$0$} (q1)
    (q5) edge[bend left=25] node[above] {$0$} (q1)
    (q5) edge[bend right=35] node[below] {$1$} (q2);
  \node[anchor=west, align=left, font=\scriptsize] at ($(q0.west)+(-0.15,-1.55)$) {
    Output: $\rho(5)=1$; $\rho(q)=0$ for $q\neq 5$.\\
    (Other mismatch edges omitted for clarity; full table used in $\Phi$.)
  };
\end{tikzpicture}
\caption{Reference recognizer $Z_{01110}$: progress states for string $01110$ (accept pulse at state $5$).}
\label{fig:z-pattern}
\end{figure}

\begingroup
\renewcommand{\arraystretch}{1.2}
\footnotesize
\begin{center}
\begin{tabularx}{0.78\linewidth}{@{}>{\raggedright\arraybackslash}p{2.8cm}Y@{}}
    \rowcolor{tableheader}
    \headingfont\bfseries Component & \headingfont\bfseries $Z_{01110}$ \\
    \addlinespace[2pt]
    $S,I,O$ & $\{0,\ldots,5\}$, $\{0,1\}$, $\{0,1\}$ \\
    \addlinespace[2pt]
    $s_0$ & $0$ \\
    \addlinespace[2pt]
    $\delta$ & Next-state table of Figure~\ref{fig:z-pattern} (overlap-aware) \\
    \addlinespace[2pt]
    $\rho(q)$ & $1$ if $q=5$, else $0$ \\
\end{tabularx}
\end{center}
\captionof{table}{Finite string-recognizer reference system.}
\label{tab:z-pattern}
\endgroup
\vspace{4pt}

\paragraph{Compiled $\Phi_{\mathrm{dyn}}$.}
Readouts in full:
\begin{align*}
  &G\bigl(\mathit{state}(5) \rightarrow \mathit{output}(1)\bigr), \\
  &G\bigl(\mathit{state}(q) \rightarrow \mathit{output}(0)\bigr)
    \quad\text{for } q\in\{0,1,2,3,4\}.
\end{align*}
Plus a guided/autonomous clause for every state and bit from the next-state table, e.g.\
$G\bigl((\mathit{state}(4)\wedge\mathit{input}(0)) \rightarrow X(\mathit{state}(5))\bigr)$.
Denote the compiled set by $\Phi_{Z_{01110}}$; it is fixed for the rest of this subsection.

\paragraph{Lifecycle iterations.}

\textit{v1 --- prototype $Z_{01110}^{\mathrm{v1}}$.}
Same automaton as Figure~\ref{fig:z-pattern}; identity maps.

\textit{v2 --- instrumented $Z_{01110}^{\mathrm{v2}}$.}
State $(q,c)$ with $q\in\{0,\ldots,5\}$ the recognizer progress and $c\in\{0,\ldots,5\}$ a parallel ones-counter modulo $6$. On each bit, $q$ updates by the reference $\delta$; $c$ increments on input $1$. Output depends only on $q$. Classical $h_S$ remains the trivial projection $(q,c)\mapsto q$.

\begingroup
\renewcommand{\arraystretch}{1.2}
\footnotesize
\begin{center}
\begin{tabularx}{0.78\linewidth}{@{}>{\raggedright\arraybackslash}p{2.8cm}Y@{}}
    \rowcolor{tableheader}
    \headingfont\bfseries Component & \headingfont\bfseries $Z_{01110}^{\mathrm{v2}}$ \\
    \addlinespace[2pt]
    $S$ & $\{0,\ldots,5\}\times\{0,\ldots,5\}$ \\
    \addlinespace[2pt]
    $s_0$ & $(0,0)$ \\
    \addlinespace[2pt]
    $\delta((q,c),b)$ & $(\delta_{01110}(q,b),\; c')$ where $c'=c{+}1\bmod 6$ if $b{=}1$, else $c$ \\
    \addlinespace[2pt]
    $\rho(q,c)$ & $\rho_{01110}(q)$ \\
    \addlinespace[2pt]
    Maps & $h_S(q,c)=q$,\ $h_I=\mathrm{id}$,\ $h_O=\mathrm{id}$ \\
\end{tabularx}
\end{center}
\captionof{table}{Instrumented recognizer (v2): parallel mod-$6$ counter is invisible to $\Phi_{Z_{01110}}$.}
\label{tab:z-pattern-med}
\endgroup
\vspace{4pt}

\textit{v3 --- shift-register refactor $Z_{01110}^{\mathrm{v3}}$.}
Here the implementation drops the explicit progress component. State is a filling-then-sliding history of the last at most five input bits (empty at start; once length $5$, each new bit drops the oldest). Output is $1$ iff the buffer equals $01110$, else $0$. There is no stored $q\in\{0,\ldots,5\}$.

\begingroup
\renewcommand{\arraystretch}{1.2}
\footnotesize
\begin{center}
\begin{tabularx}{0.78\linewidth}{@{}>{\raggedright\arraybackslash}p{2.8cm}Y@{}}
    \rowcolor{tableheader}
    \headingfont\bfseries Component & \headingfont\bfseries $Z_{01110}^{\mathrm{v3}}$ \\
    \addlinespace[2pt]
    $S$ & $\{ w\in\{0,1\}^{\le 5}\}$ (histories of length at most $5$) \\
    \addlinespace[2pt]
    $s_0$ & $\varepsilon$ (empty) \\
    \addlinespace[2pt]
    $\delta(w,b)$ & $wb$ if $|w|<5$; else $\mathrm{tail}(w)\,b$ \\
    \addlinespace[2pt]
    $\rho(w)$ & $1$ iff $w=01110$; else $0$ \\
    \addlinespace[2pt]
    Maps & $h_S(w)=\delta^{*}_{01110}(0,w)$ (fold the reference table over $w$) \\
\end{tabularx}
\end{center}
\captionof{table}{Shift-register recognizer (v3): no explicit progress state; $h_S$ is non-trivial.}
\label{tab:z-pattern-cplx}
\endgroup
\vspace{4pt}

\paragraph{Workflow contrast at v3.}
\textit{Assertional:} $\Phi_{Z_{01110}}$ is still the fragment compiled at the reference; verifying the refactor means re-checking that same clause set on the shift-register tables (on these finite spaces, equivalently searching or verifying a Def.~4.3 witness). No new property derivation is required from the engineer.
\textit{Classical:} the engineer must construct and prove a map $h_S$ from every history $w$ to a progress state in $\{0,\ldots,5\}$. The correct choice is to run the reference transition table from $0$ along $w$ (equivalently, longest-prefix-suffix progress of $01110$ on $w$). Surjectivity holds because each progress $q$ is realized by the length-$q$ prefix of $01110$. Transition preservation is the statement that sliding the window agrees with one reference step---true for this overlap-aware table, but not obvious a priori and not a component projection. Readout preservation is $\rho_{\mathrm{v3}}(w)=1$ iff $h_S(w)=5$. Every subsequent architectural change of this kind again forces a human to rebuild and re-prove such an $h_S$, whereas the assertional side reuses $\Phi_{Z_{01110}}$.

Instrumentation at v2 keeps a trivial projection; the v3 shift register is where classical maintenance becomes laborious while the assertional artifact stays fixed. Graphically, Figure~\ref{fig:z-pattern} remains the required behavior; checking $\Phi_{Z_{01110}}$ asks whether the buildable's projected progress obeys that table.

\subsection{Dual-Port Pattern Select}
\label{sec:case-dual}

A light extension shows mode-dependent required behavior. Input is a pair $(\mathrm{mode},\mathrm{data})$ with $\mathrm{mode}\in\{\mathrm{short},\mathrm{long}\}$ and $\mathrm{data}\in\{0,1\}$. Mode selects which pattern the shared progress state is matching: short mode tracks $01$; long mode tracks $01110$. Output is an accept pulse when the selected pattern has just completed.

\begingroup
\renewcommand{\arraystretch}{1.2}
\footnotesize
\begin{center}
\begin{tabularx}{0.78\linewidth}{@{}>{\raggedright\arraybackslash}p{2.8cm}Y@{}}
    \rowcolor{tableheader}
    \headingfont\bfseries Component & \headingfont\bfseries $Z_{\mathrm{dual}}$ (sketch) \\
    \addlinespace[2pt]
    $I$ & $\{\mathrm{short},\mathrm{long}\}\times\{0,1\}$ \\
    \addlinespace[2pt]
    $S$ & Progress states for the active pattern (finite) \\
    \addlinespace[2pt]
    $\delta$ & Advance the short or long automaton according to mode \\
    \addlinespace[2pt]
    $\rho$ & $1$ on accept of the active pattern; else $0$ \\
\end{tabularx}
\end{center}
\captionof{table}{Dual-port select: mode chooses which recognizer dynamics apply.}
\label{tab:z-dual}
\endgroup
\vspace{4pt}

Compiled $\Phi_{\mathrm{dyn}}$ is again a finite clause list from that reference table. Direct and elaborated buildables (extra logging beside the progress component) are checked against it exactly as in the ones-counter and single-pattern cases above.
